\chapter{两个剩余位置：按输入选择返回}\label{ch:two}

\section{为什么把两个剩余位置单列}
本讲取 $r=2$、$n=2p+2$，使用
\[
(A^p,B^p,C^2)=((u+v)^p,(u-v)^p,(-pu)^2).
\]
令 $m=p+1$，合法性化为 $v$ 与 $m$ 互素，尤其 $v$ 为奇数。
这个最紧的偶数余量不走共同终端中生成全部有限单位的路线，
而是根据当前输入选择两次数字返回，直接把第二坐标改成 $\pm1$ 模一个明确的深层。

中国剩余定理在这里的角色很具体：第一次使第二坐标在控制层处处为单位，
第二次把它的完整余数指定为 $\pm1$。两次都同时安排本原性，
避免最后约去公因子时破坏刚设好的余数。

\begin{proposition}[余量为2的核心]\label{prop:r2}
若 $p\ge5$，全部合法 $K_{p,2}$ 可归零。
\end{proposition}
\begin{proof}
引理\ref{lem:two-carrier}给出 $U(\Z[1/p])$ 和一个明确整数深核。其两次 CRT 数字返回使第二坐标等于 $\pm1$ 模该深核层，再由引理\ref{lem:fiber}精确送入相反值终端。
\end{proof}



\section{完整构造}
为了方便跟踪，以下把构造分为四个阶段：
上平移与缩放的逆、横向反射与深层、第一次数字返回、第二次数字返回。

\begin{lemma}[两次CRT将尺度送到正负1]\label{lem:two-carrier}
在 $((u+v)^p,(u-v)^p,(-pu)^2)$ 上，合法本原参数满足 $\gcd(u,v)=\gcd(v,p+1)=1$。全部合法参数可送到 $(u,v)=(0,1)$ 的射影方向。
\end{lemma}
\begin{proof}

\par\noindent\textbf{第一阶段：上平移与缩放的正向逆。}
记 $m=p+1$。保留两份 $A=u+v$，平均
\[(A^{(p-3)/2},B^{(p+1)/2},C),\qquad
(A^{(p-1)/2},B^{(p-1)/2},C)
\]
实现 $-U(1)/p$；换块给 $U(-1)$。互补平衡混合给 $F=\diag(1,1/p)$。其逆由四次平均实现：先平均两组 $(A^{(p-1)/2},B^{(p-1)/2},C)$，产生 $2p$ 份 $D=-u/p$；保留两份 $D$，分别平均一份剩余 $A$ 或 $B$ 与 $p-1$ 份 $D$。输出为 $(u/p^2,v/p)$。因此 $U(\Z[1/p])$ 可实现。

\par\noindent\textbf{第二阶段：三步反射与深同余群。}
再平均 $(A^{(p+1)/2},B^{(p-3)/2},C)$。保留两份新值，按 $(x,x+1,b,(p-3)/2)$ 选下一组，其中 $b=1$（$p\equiv1\pmod4$）或 $0$（$p\equiv3\pmod4$），$x=(p-b-(p-1)/2)/2$，再平均补组。返回为 $J/p^2$，
\[
J=\begin{pmatrix}1&-2\\c&-1\end{pmatrix},\qquad
c=\begin{cases}-(p^2-1)/2,&p\equiv1\pmod4,\\(p+1)^2/2,&p\equiv3\pmod4.
\end{cases}
\]
两类计数分别是 $(k,k+1,1,2k-1)$ 和 $(k+1,k+2,0,2k)$，均在库存内。$J^2=(1-2c)I$。若 $S=\diag(1,-1)$、$H=JU(2)$，则 $[S,H]=L(-2c)$。借 $F$ 共轭得到完整 $L(2cR)$，$R=\Z[1/p]$。令 $q=(2c)^2$，引理\ref{lem:deep}给 $\Gamma(q,R)$；$m\mid2c$ 且 $\gcd(p,q)=1$。


\par\noindent\textbf{第三阶段：使第二坐标在整个控制层可逆。}
对 $2\le b\le p$，保留两份 $A$，平均 $(A^{b-2},B^{p-b},C^2)$ 和 $(A^{p-b},B^b)$，实现
\[
p^{-1}M_b,\qquad M_b=\begin{pmatrix}-1&-1\\-m&2b-m\end{pmatrix}.
\]

这个 $M_b$ 是可用的正向返回，但不预先假定它对所有本原输入安全或可逆。
下面每次使用它，都会先安排输入，使本原输出保持合法。
由于 $v$ 奇，先用整数上平移把 $u$ 改成 $2a-v$，即选 $2a\equiv u\pmod v$。
直接乘矩阵可见
\[
\frac1pM_p(2a-v,v)^T=\frac2p(-a,pv-ma)^T.
\]
其中 $2/p$ 是全部核心的共同尺度，所余整数列将被安排为本原。
执行 $M_p$ 后本原分析列为 $(-a,pv-ma)$。对 $\ell\mid q$，若 $\ell\mid m$，第二项自动非零；若 $\ell\mid v$，$a$ 自动非零；其余素数指定 $a=(pv-1)/m\pmod\ell$。
这些额外同余只在不整除 $v$ 的素数处新增，所以与 $2a\equiv u\pmod v$ 相容。
再令 $p\nmid a$；若 $p\mid v$，原来的同余已经保证这一点，否则另加 $a\equiv1\pmod p$。
由于 $\gcd(a,v)=1$，有
\[
\gcd(a,pv-ma)=\gcd(a,p)=1.
\]
这给出本原性，而对所有 $\ell\mid q$ 的选择给出第二坐标的单位性。
CRT 给出本原输出 $(u_1,v_1)$ 且 $\gcd(v_1,q)=1$。

\par\noindent\textbf{第四阶段：直接指定第二坐标的完整余数。}
取 $v_1^{-1}\pmod m$ 的代表 $b_0\in[1,p]$；若 $b_0=1$，改取 $b=p,\epsilon=-1$，否则取 $b=b_0,\epsilon=1$。于是 $b$ 奇、$2\le b\le p$、$bv_1\equiv\epsilon\pmod m$。选
\[
2a_1\equiv u_1\pmod{v_1},\qquad
a_1\equiv(bv_1-\epsilon)/m\pmod{q/m},\qquad \gcd(a_1,b)=1.
\]

除法 $(bv_1-\epsilon)/m$ 是整数除法，因为先前已选 $bv_1\equiv\epsilon\pmod m$。
此时第二个同余恰好等价于 $bv_1-ma_1\equiv\epsilon\pmod q$。
前两个模数互素；$b$ 的素因子若整除其中一个模数，相应条件已使 $a_1$ 非零，其余处由 CRT 补齐。平移后执行 $M_b$，得到本原列 $(-a_1,bv_1-ma_1)$，第二项为 $\epsilon\pmod q$。
因为 $\gcd(a_1,v_1)=\gcd(a_1,b)=1$，有
\[
\gcd(a_1,bv_1-ma_1)=\gcd(a_1,bv_1)=1.
\]
所以这里无须再约去任何内容，已经安排好的第二坐标余数确实保留。
再用上平移使第一项为零模 $q$，由引理\ref{lem:fiber}精确到 $(0,\epsilon)$。相反值收尾完成。
\end{proof}

